\documentclass[9pt,a4paper]{extarticle}

% ======================
% Codificação, idioma e layout
% ======================
\usepackage[T1]{fontenc}
\usepackage[utf8]{inputenc}
\usepackage[brazil]{babel}
\usepackage{geometry}
\geometry{margin=2cm}
\usepackage{parskip}
\usepackage{setspace}
\usepackage{microtype}

% ======================
% Matemática e símbolos
% ======================
\usepackage{amsmath, amssymb, amsthm, bm}

% ======================
% Figuras e tabelas
% ======================
\usepackage{graphicx}
\usepackage{booktabs}
\usepackage{array}
\usepackage{tabularx}
\usepackage{float}
\usepackage{caption}
\usepackage{subcaption}

% ======================
% Hiperlinks
% ======================
\usepackage{hyperref}
\hypersetup{
  colorlinks=true,
  linkcolor=blue,
  urlcolor=blue
}

% ======================
% Início do documento
% ======================
\begin{document}

\begin{center}
  \Large \textbf{Estudo Dirigido – Condições de Primeira e Segunda Ordem para Problemas com Restrições} \\[4pt]
  \normalsize Disciplina: Otimização Não Linear
\end{center}

\begin{enumerate}


% ============================================================
% 5. Interpretação das KKT
% ============================================================
\item \textbf{Interpretação das Condições KKT}

Considere o problema geral:
\[
\min_{\mathbf{x}} \; f(\mathbf{x})
\quad\text{sujeito a}\quad 
g_i(\mathbf{x}) = 0,\; i=1,\ldots,m,\qquad
h_j(\mathbf{x}) \le 0,\; j=1,\ldots,p.
\]

As Condições de Karush--Kuhn--Tucker (KKT) de primeira ordem são:
\[
\begin{cases}
\text{(1) Estacionariedade:} 
& \nabla f(\mathbf{x}^*) 
+ \displaystyle\sum_{i=1}^m \lambda_i^* \nabla g_i(\mathbf{x}^*)
+ \displaystyle\sum_{j=1}^p \mu_j^* \nabla h_j(\mathbf{x}^*)
= 0, \\[6pt]

\text{(2) Viabilidade primária:}
& g_i(\mathbf{x}^*) = 0,\quad i=1,\ldots,m, \\
& h_j(\mathbf{x}^*) \le 0,\quad j=1,\ldots,p, \\[6pt]

\text{(3) Viabilidade dual:} 
& \mu_j^* \ge 0,\quad j=1,\ldots,p, \\[6pt]

\text{(4) Complementaridade:}
& \mu_j^* \, h_j(\mathbf{x}^*) = 0,\quad j=1,\ldots,p.
\end{cases}
\]

Responda:

\begin{enumerate}
    \item Nas consições apresentadas em aula, não colocamos $\mu_j^* \geq 0$. Utilize um problema com apenas uma restrição de desigualdade para explicar a necessidade desta condição.  
    \item Explique detalhadamente o significado de cada uma das três condições restantes, incluindo:
    \begin{enumerate}
        \item a interpretação geométrica de cada condição;
        \item o papel de cada condição na eliminação de direções de descida factíveis;
        \item o papel dos multiplicadores na interpretação das restrições;
        \item por que a complementaridade distingue quais restrições influenciam o ótimo.
    \end{enumerate}
\end{enumerate}




% ============================================================
% 1. Direções factíveis e de melhora
% ============================================================
\item \textbf{Direções factíveis e de melhora}

Considere o problema:
\[
\min f(x,y)=x^2+y^2 
\qquad \text{sujeito a } 
g(x,y)=x+y-1=0.
\]

\begin{enumerate}
    \item Encontre todas as direções factíveis no ponto $(x,y)=(0,1)$.
    \item Determine se existe alguma direção de melhora nesse ponto.
    \item Conclua se $(0,1)$ pode ser um ponto ótimo local.
\end{enumerate}


% ============================================================
% 2. Lagrangeana e KKT básicas
% ============================================================
\item \textbf{Construção da Lagrangeana e condições de primeira ordem}

Considere o problema:
\[
\min_{x,y} \; f(x,y) = x^2 + y^2 - 4x
\]
\[
\text{sujeito a } h(x,y)=x+y-2\le 0.
\]

\begin{enumerate}
    \item Escreva a Lagrangeana $\mathcal{L}(x,y,\mu)$.
    \item Encontre as condições KKT de primeira ordem.
    \item Determine os candidatos a ponto ótimo resolvendo o sistema KKT.
    \item Classifique a restrição como ativa ou inativa no ponto obtido.
\end{enumerate}


% ============================================================
% 3. KKT com múltiplas restrições
% ============================================================
\item \textbf{KKT com múltiplas restrições e interpretação geométrica}

Considere o problema:
\[
\min_{x,y} \; f(x,y) = x 
\]
\[
\text{sujeito a }
\begin{cases}
g_1(x,y)=x^2+y^2-1=0, \\
g_2(x,y)=y- x \le 0.
\end{cases}
\]

\begin{enumerate}
    \item Desenhe qualitativamente a região viável.
    \item Liste quais restrições são ativas no ponto que minimiza $x$.
    \item Escreva as condições KKT para esse problema.
    \item Mostre que o ponto ótimo deve satisfazer simultaneamente $g_1(x,y)=0$ e $g_2(x,y)=0$.
    \item Encontre o ponto ótimo e os multiplicadores correspondentes.
\end{enumerate}


% ============================================================
% 4. Segunda ordem com igualdade
% ============================================================
\item \textbf{Condições de segunda ordem com restrições de igualdade}

As condições de segunda ordem analisam a curvatura da Lagrangeana ao longo das direções factíveis.  
Para um problema geral
\[
\min_{\bm{x}} f(\bm{x})
\quad \text{sujeito a} \quad g_i(\bm{x}) = 0,\; i=1,\dots,m,
\]
se $(\bm{x}^*, \bm{\lambda}^*)$ satisfaz as condições de primeira ordem, define-se o Hessiano da Lagrangeana em relação a $\bm{x}$:
\[
\nabla^2_{\bm{x}\bm{x}} \mathcal{L}(\bm{x}^*, \bm{\lambda}^*)
= \nabla^2 f(\bm{x}^*) + \sum_{i=1}^m \lambda_i^* \nabla^2 g_i(\bm{x}^*).
\]
Seja ainda o espaço das direções tangentes factíveis no ponto ótimo:
\[
\mathcal{D} = \{\bm{d} \neq 0 : \nabla g_i(\bm{x}^*)^\top \bm{d} = 0,\; i=1,\dots,m\}.
\]

Então:
\begin{itemize}
    \item \textbf{Condição necessária de segunda ordem:}
    \[
    \bm{d}^\top \nabla^2_{\bm{x}\bm{x}} \mathcal{L}(\bm{x}^*, \bm{\lambda}^*) \bm{d} \ge 0,
    \quad \forall \bm{d} \in \mathcal{D}.
    \]
    \item \textbf{Condição suficiente (mínimo local):}
    \[
    \bm{d}^\top \nabla^2_{\bm{x}\bm{x}} \mathcal{L}(\bm{x}^*, \bm{\lambda}^*) \bm{d} > 0,
    \quad \forall \bm{d} \in \mathcal{D}.
    \]
\end{itemize}

Considere agora o problema específico:
\[
\min f(x,y)= x^2 +xy + y^2
\]
\[
\text{sujeito a } g(x,y)=x-y=0.
\]

\begin{enumerate}
    \item Escreva a Lagrangeana e encontre $(x^*,y^*,\lambda^*)$.
    \item Encontre o espaço das direções tangente factíveis no ponto ótimo.
    \item Calcule o Hessiano da Lagrangeana $\mathcal{L}_{xx}$.
    \item Verifique a condição necessária de segunda ordem.
    \item Verifique a condição suficiente de segunda ordem.
    \item Classifique o ponto como mínimo, máximo ou ponto de sela dentro da variedade factível.
\end{enumerate}

\end{enumerate}



\end{document}

